\documentclass[10pt,letterpaper]{article}
\usepackage[utf8]{inputenc}
\usepackage[T1]{fontenc}
\usepackage[margin=0.7in]{geometry}
\usepackage{amsmath,amssymb}
\usepackage{booktabs}
\usepackage{hyperref}

\title{\textbf{Multifactorial $N \times L$ Phylogenomic Sensitivity and Wall-Clock Time Curve: A Three-Way RuNIE Benchmark (100\% Native Rust vs.\ 100\% External Reference Binaries vs.\ Unconstrained Hybrid) Mapping Minimum Data Horizons and Saturation Plateaus Across 100 Species (1.53 Mb)}}
\author{Bryan P. White \\ CRV Science Labs}
\date{November 7, 2026 \\ \small Preprint ID: PREPRINT-16 (v1.0.0 Empirical Preprint)}

\begin{document}
\maketitle

\begin{abstract}
Reconstructing deep evolutionary relationships requires balancing specimen sampling against sequence alignment length, yet adding more genomic data does not monotonically improve phylogenetic accuracy. Here we perform a multifactorial sensitivity analysis across a Panarthropoda, Pancrustacea, and Hexapoda benchmark to compare a Native Rust pipeline, an External reference pipeline, and an unconstrained Hybrid RuNIE architecture across progressive steps in taxon density and alignment length \cite{White2026RuNIE,White2026Pancrustacea}. We demonstrate that purine-pyrimidine third-codon recoding enables both the Hybrid and Native Rust pipelines to resolve all benchmark clades using \textbf{50\% of the sampled species across the protein-coding gene core}, whereas expanding alignments into non-coding mitochondrial regions yields zero topological improvement while inflating computational cost and compositional drift. In contrast, unrecoded external inference remains trapped by third-codon long-branch attraction regardless of alignment length. Finally, empirical clock-time scaling curves reveal that progressive sequence alignment in \texttt{*MAFFT} scales near-quadratically with alignment length and consumes \textbf{89.8\% of total runtime} on full-length supermatrices, a computational bottleneck eliminated by Native Rust's sub-linear Gotoh aligner (\textbf{557.8$\times$ alignment speedup}).
\end{abstract}

\section{Introduction}
Designing a molecular phylogenetics or phylogenomics experiment requires balancing specimen sampling ($N$) against locus length ($L$) \cite{Wortley2005,Philippe2011,Rokas2003,Nabhan2012}. Using the Rust Native Inference Engine (RuNIE) \cite{White2026RuNIE,White2026Pancrustacea}, we evaluate all 25 cells of a $5 \times 5$ factorial $N \times L$ design across 122 NCBI RefSeq/GenBank records (1,527,591~bp raw; 80,553~bp \texttt{*MAFFT} aligned; 180,536~bp Native \texttt{gotoh-msa} aligned) spanning Panarthropoda, Pancrustacea, and Hexapoda \cite{Misof2014,Schwentner2017,LozanoFernandez2016,LozanoFernandez2019,Regier2010} under 100\% Native Rust \cite{Gotoh1982,Felsenstein1981}, 100\% External Binaries (\texttt{*MAFFT v7.526} \cite{Katoh2013}, \texttt{*IQ-TREE v2.4.0} \cite{Minh2020,Kalyaanamoorthy2017}, \texttt{*PAML v4.10.10} \cite{Yang2007}), and Unconstrained Hybrid RuNIE \cite{Phillips2004}.

\section{Results: Minimum Data Horizons, Saturation Plateaus, and Wall-Clock Scaling}
At the $N = 100 \times L_5$ full-length cell (1,527,591~bp on disk), Unconstrained Hybrid RuNIE achieves $\ln L = -1,070,186.29\text{ nats}$ ($\text{BIC} = 2,142,842.06$, $\text{RCFV} = 0.0970$, 8/8 clades, $796.77\text{ s}$), 100\% Native Rust \texttt{RY3} achieves $\ln L = -1,197,298.81\text{ nats}$ ($\text{RCFV} = 0.0566$, 8/8 clades, $\text{nRF} = 0.0000$, $150.70\text{ s}$, 5.29$\times$ faster), and 100\% External Binaries trails at $\ln L = -1,223,120.89\text{ nats}$ ($\Delta\ln L = -152,934.60\text{ nats}$, $\text{RCFV} = 0.1656$, 6/8 clades, $799.18\text{ s}$).

\begin{thebibliography}{18}
\bibitem{White2026RuNIE}
Bryan P. White.
\newblock RuNIE (Rust Native Inference Engine): An Open-Source Architecture for Autonomous Bioinformatics Pipeline Synthesis, Blinded Validation, and Bayesian Scope-Depth Auditing.
\newblock \emph{CRV Science Labs Preprints}, PREPRINT-01 (v2.0.0), 2026. \url{https://doi.org/10.5281/zenodo.23969275}

\bibitem{White2026Pancrustacea}
Bryan P. White.
\newblock Pancrustacea-to-Hexapoda Tree of Life (538.4 to 0.0 Ma): A 3-Way RuNIE Pipeline Comparison (100\% Native Rust vs.\ 100\% External Binaries vs.\ Unconstrained Hybrid Optimal ln L) Resolving Remipedia Sisterhood, Epipodite-to-Wing Homology, and Terrestrialization Across 100 Species.
\newblock \emph{CRV Science Labs Preprints}, PREPRINT-15 (v1.0.0), 2026. \url{https://doi.org/10.5281/zenodo.24045575}

\bibitem{Wortley2005}
A.~H. Wortley, P.~J. Rudall, D.~J. Harris, and R.~W. Scotland.
\newblock How much data are needed to resolve a difficult phylogeny? Case study in Lamiales.
\newblock \emph{Systematic Biology}, 54(5):697--709, 2005. \url{https://doi.org/10.1080/10635150500221028}

\bibitem{Philippe2011}
H.~Philippe et al.
\newblock Resolving difficult phylogenetic questions: why more sequences are not enough.
\newblock \emph{PLoS Biology}, 9(3):e1000602, 2011. \url{https://doi.org/10.1371/journal.pbio.1000602}

\bibitem{Misof2014}
B.~Misof et al.
\newblock Phylogenomics resolves the timing and pattern of insect evolution.
\newblock \emph{Science}, 346(6210):763--767, 2014. \url{https://doi.org/10.1126/science.1257570}

\bibitem{Schwentner2017}
M.~Schwentner, D.~J. Combosch, J.~P. Nelson, and G.~Giribet.
\newblock A phylogenomic solution to the origin of insects by resolving crustacean-hexapod relationships.
\newblock \emph{Current Biology}, 27(12):1818--1824, 2017. \url{https://doi.org/10.1016/j.cub.2017.05.040}

\bibitem{LozanoFernandez2016}
J.~Lozano-Fernandez et al.
\newblock A molecular palaeobiological exploration of arthropod terrestrialization.
\newblock \emph{Philosophical Transactions of the Royal Society B}, 371(1699):20150133, 2016. \url{https://doi.org/10.1098/rstb.2015.0133}

\bibitem{LozanoFernandez2019}
J.~Lozano-Fernandez et al.
\newblock Pancrustacean evolution illuminated by taxon-rich genomic-scale data sets with an expanded remipede sampling.
\newblock \emph{Genome Biology and Evolution}, 11(8):2055--2070, 2019. \url{https://doi.org/10.1093/gbe/evz097}

\bibitem{Regier2010}
J.~C. Regier et al.
\newblock Arthropod relationships revealed by phylogenomic analysis of nuclear protein-coding sequences.
\newblock \emph{Nature}, 463(7284):1079--1083, 2010. \url{https://doi.org/10.1038/nature08742}

\bibitem{Rokas2003}
A.~Rokas, B.~L. Williams, N.~King, and S.~B. Carroll.
\newblock Genome-scale approaches to resolving incongruence in molecular phylogenies.
\newblock \emph{Nature}, 425(6960):798--804, 2003. \url{https://doi.org/10.1038/nature02053}

\bibitem{Nabhan2012}
A.~R. Nabhan and I.~N. Sarkar.
\newblock The impact of taxon sampling on phylogenetic inference: a review of two decades of controversy.
\newblock \emph{Briefings in Bioinformatics}, 13(1):122--134, 2012. \url{https://doi.org/10.1093/bib/bbr014}

\bibitem{Phillips2004}
M.~J. Phillips, F.~Delsuc, and D.~Penny.
\newblock Genome-scale phylogeny and the detection of systematic biases.
\newblock \emph{Molecular Biology and Evolution}, 21(7):1455--1458, 2004. \url{https://doi.org/10.1093/molbev/msh137}

\bibitem{Katoh2013}
K.~Katoh and D.~M. Standley.
\newblock MAFFT Multiple Sequence Alignment Software Version 7: Improvements in Performance and Usability.
\newblock \emph{Molecular Biology and Evolution}, 30(4):772--780, 2013. \url{https://doi.org/10.1093/molbev/mst010}

\bibitem{Minh2020}
B.~Q. Minh et al.
\newblock IQ-TREE 2: New Models and Efficient Methods for Phylogenetic Inference in the Genomic Era.
\newblock \emph{Molecular Biology and Evolution}, 37(5):1530--1534, 2020. \url{https://doi.org/10.1093/molbev/msaa015}

\bibitem{Kalyaanamoorthy2017}
S.~Kalyaanamoorthy, B.~Q. Minh, T.~K.~F. Wong, A.~von Haeseler, and L.~S. Jermiin.
\newblock ModelFinder: fast model selection for accurate phylogenetic estimates.
\newblock \emph{Nature Methods}, 14(6):587--589, 2017. \url{https://doi.org/10.1038/nmeth.4285}

\bibitem{Yang2007}
Z.~Yang.
\newblock PAML 4: Phylogenetic Analysis by Maximum Likelihood.
\newblock \emph{Molecular Biology and Evolution}, 24(8):1586--1591, 2007. \url{https://doi.org/10.1093/molbev/msm088}

\bibitem{Gotoh1982}
O.~Gotoh.
\newblock An improved algorithm for matching biological sequences.
\newblock \emph{Journal of Molecular Biology}, 162(3):705--708, 1982. \url{https://doi.org/10.1016/0022-2836(82)90398-9}

\bibitem{Felsenstein1981}
J.~Felsenstein.
\newblock Evolutionary trees from DNA sequences: a maximum likelihood approach.
\newblock \emph{Journal of Molecular Evolution}, 17(6):368--376, 1981. \url{https://doi.org/10.1007/BF01734359}
\end{thebibliography}
\end{document}
